\begin{table}[H]
\centering\centering\centering
\caption{Roster racial composition and game outcomes: within-season design, regular season 2002-2025 \label{tab:roster-diversity-game}}
\centering
\resizebox{\ifdim\width>\linewidth\linewidth\else\width\fi}{!}{
\begin{tabular}[t]{lccccccccc}
\toprule
  & (1) ATS margin & (2) ATS, residual share & (3) Margin, spread and home & (4) Win minus market prob. & (5) ATS, franchise and season FE & (6) ATS, opponent-season FE & (7) Margin, opponent-season FE & (8) Symmetry check & (9) Symmetry, game QB\\
\midrule
Active-roster share Black & -2.293 &  & -1.170 & -0.041 & -6.429** & -5.077 & -0.617 & -5.161 & -4.734\\
 & (8.425) &  & (8.331) & (0.278) & (2.680) & (9.383) & (8.509) & (11.744) & (11.420)\\
Residual active-roster share Black &  & -6.047 &  &  &  &  &  &  & \\
 &  & (8.428) &  &  &  &  &  &  & \\
Opponent active-roster share Black &  &  &  &  &  &  &  & 11.279 & 12.094\\
 &  &  &  &  &  &  &  & (12.573) & (12.395)\\
Point spread (team favored $>$ 0) &  &  & 0.801*** &  &  &  & -0.202*** &  & \\
 &  &  & (0.020) &  &  &  & (0.041) &  & \\
Home team &  &  & 0.947*** &  &  &  & 5.243*** &  & \\
 &  &  & (0.235) &  &  &  & (0.288) &  & \\
Starting QB Black (game) &  &  &  &  &  &  &  &  & -1.037\\
 &  &  &  &  &  &  &  &  & (1.765)\\
Opponent starting QB Black (game) &  &  &  &  &  &  &  &  & -0.893\\
 &  &  &  &  &  &  &  &  & (1.995)\\
\midrule
Observations & 12446 & 12446 & 12446 & 12446 & 12446 & 12446 & 12446 & 6223 & 6223\\
R$^2$ & 0.065 & 0.065 & 0.260 & 0.063 & 0.008 & 0.128 & 0.360 & 0.250 & 0.250\\
Within R$^2$ & 0.001 & 0.002 & 0.102 & 0.002 & 0.002 & 0.002 & 0.036 & 0.006 & 0.007\\
Two-way clustered SE, roster share & (8.458) & (8.493) & (8.353) & (0.281) & (2.703) & (9.414) & (8.520) & (9.518) & (8.853)\\
Two-way clusters & Franchise, game & Franchise, game & Franchise, game & Franchise, game & Franchise, game & Franchise, game & Franchise, game & Home, away team & Home, away team\\
WCB p-value, roster share & 0.786 & 0.481 & 0.888 & 0.889 & 0.023 &  &  &  & \\
p-value, own + opponent share = 0 &  &  &  &  &  &  &  & 0.683 & 0.625\\
Mean of outcome & 0.000 & 0.000 & 0.000 & 0.000 & 0.000 & 0.000 & 0.000 & 0.049 & 0.049\\
SD of roster share net of FE & 0.017 & 0.015 & 0.017 & 0.017 & 0.043 & 0.016 & 0.016 & 0.015 & 0.015\\
Effect of 1-SD change in expected share & -0.039 & -0.093 & -0.020 & -0.001 & -0.279 & -0.083 & -0.010 & -0.079 & -0.072\\
Sample & Both teams & Both teams & Both teams & Both teams & Both teams & Both teams & Both teams & One row per game & One row per game\\
Franchise $\times$ season FE & Yes & Yes & Yes & Yes & No & Yes & Yes & Yes & Yes\\
Franchise and season FE & No & No & No & No & Yes & No & No & No & No\\
Opponent $\times$ season FE & No & No & No & No & No & Yes & Yes & Yes & Yes\\
Mean prior and prior-covariate shares & Yes & Yes & Yes & Yes & Yes & Yes & Yes & Yes & Yes\\
\bottomrule
\multicolumn{10}{l}{\rule{0pt}{1em}* p $<$ 0.1, ** p $<$ 0.05, *** p $<$ 0.01}\\
\end{tabular}}
\end{table}

{\footnotesize
\noindent\textit{Notes:} This table includes the estimation results of equation (3), the within-season game-level version of equation (2). The unit of observation is a franchise-game in the regular season, 2002-2025 (12,446 team-games; each game enters once for each team, except in columns (8)-(9)). The treatment is the Black share of the game-day active roster (players not listed inactive), weighting players equally. Inactive designations are not recorded in 2016-2018, when the active roster equals the full game-day roster. With franchise $\times$ season fixed effects, identification comes from week-to-week changes in a team's active roster (injuries, inactives, signings) within a team-season. Under the predicted measure every column includes the calibration controls of the active roster (its mean prior P(Black) and the weighted shares of the prior's draft-round, college-type, rookie-era and hometown-county levels); columns (8)-(9) also include the opponent's, and column (9) the quarterbacks' priors. Column (1), the main specification, uses the margin against the closing spread (margin minus the spread, positive when the team is favored) as the outcome, which imposes a unit slope on the spread; the raw slope of the margin on the spread is about 1.04. The spread prices opponent strength and home field, so no opponent fixed effects are needed. Because the betting market prices roster changes, including composition, the coefficient estimates only the part of any composition effect the market fails to price; it is zero under market efficiency even if composition has a causal effect. Column (2) uses the residual share, net of the team's game-day position mix. The residual share is the actual share minus the expected share given the team's position mix (the sum over position groups of the team's position weight times the league-season Black share at that position on all other franchises). Column (3) uses the raw margin and controls for the spread and a home indicator (home varies within a team-season and is not absorbed). The spread's slope there is estimated within team-season, where the spread also responds to earlier results of the same team-season; it is therefore not strictly exogenous, and the slope is far below one. Column (4) is a linear probability model of a win (ties count one half) net of the market's pre-game win probability. Columns (5)-(7) show the sensitivity to the fixed effects: franchise and season fixed effects (column (5)), and franchise $\times$ season plus opponent $\times$ season fixed effects (columns (6)-(7), the plan's original design). In column (7) the spread's slope is distorted further by the two sets of incidental team-season parameters, so the column is a robustness check only. Columns (8)-(9) are a symmetry check, not a placebo: the opponent's composition can legitimately affect the margin, and in a symmetric model it should enter with the opposite sign and similar size. In the stacked sample the opponent coefficient is mechanically minus the own coefficient, so these columns use one row per game (the home team; 6,223 games), with  home-team $\times$ season and away-team $\times$ season fixed effects. The p-value row tests that the own and opponent coefficients sum to zero. Column (9) adds the race of each team's starting quarterback in that game (zero-filled with missing indicators). Standard errors, in parentheses, are clustered at the franchise level (32 clusters). The rows below report standard errors clustered two ways: by franchise and game in the stacked columns, where the two rows of a game are mirror images and share a game cluster but neither a franchise nor an opponent cluster; and by home and away franchise in columns (8)-(9). Wild cluster bootstrap-t p-values (franchise clusters, Webb weights, 9,999 replications) are  reported for columns (1)-(5); in columns (1)-(4) the team-season fixed effects are projected out before the bootstrap (team-seasons are nested in franchise clusters, so this is equivalent to including them). The SD row is the SD of the treatment net of the column's fixed effects; the effect row multiplies it by the coefficient, in points (columns (1)-(3) and (5)-(9)) or win probability (column (4)). Under the model-only predicted measure a share Black is the expected share of players who are non-Hispanic Black alone (multiracial and Hispanic Black players count as non-Black); the documented, provisional and hand-coded sensitivity measures count Black alone or in combination. Under regression calibration the coefficient on an expected roster share (the mean over players of the predicted probability of being non-Hispanic Black) equals the effect of the true share only if the probabilities are calibrated at the team level given the regression's controls; this requires that players' names and hometown counties be unrelated to team performance given race and the prior's covariates, and the calibration controls (the team mean of the prior and the weighted shares of its draft-round, college-type, rookie-era and hometown-county levels, plus the priors of the coaches, head coach and quarterback whose predicted race enters as a control) condition on those covariates. Shrinkage toward the prior by itself causes no attenuation (the error of an expected share is of the Berkson type), but individual-level calibration does not imply team-level calibration: if the true share rises more (less) than one-for-one with the expected share across teams, the coefficient overstates (understates) the effect, so the sign of any bias is not known. A league-wide level error, such as the model's shortfall against the TIDES league share, is absorbed by the season fixed effects. A team-level calibration check among documented players (validation only): regressing a franchise-season's documented Black share of its documented players on the same players' mean predicted probability, with franchise and season fixed effects, gives a slope of 0.79 (SE 0.06) for snap weights and  0.69 (SE 0.06) for headcount weights, against one under calibration. Documented players are a selected, mostly Black and famous subset (about 28 percent of the weight), so the check does not establish the team-level calibration of the full roster in either direction. Effect rows refer to a one-SD change in the expected share, whose SD is smaller than that of the true share. Race is predicted, not observed: each person's probability of being non-Hispanic Black alone combines first name, surname and hometown county (BIFSG) with an NFL-specific prior estimated by EM on predetermined characteristics (players: position at entry, rookie era, draft round, college type, county availability; staff: role, unit and era at first appearance); documented race statements are not used (notes/race-prediction-design.md). Black Hispanic and multiracial Black persons are in the other-race probability, so P(Black) is narrower than the Black-alone-or-in-combination definition of the hand codes and of TIDES. Team shares are expected shares (mean member probability). Their error is Berkson-type: estimates are consistent if the probabilities are calibrated given the regression's controls (the team mean of the prior is one of them), but less precise than with observed race; individual coaches' probabilities are noisy because names carry little information for many Black coaches. The validation exhibit compares the predictions with published TIDES shares.
\par}

